% ---------------------------------------------------------------------------
% Reference-Anchored Turn Taking (RATT)
%
% Standalone section intended to be merged into the main AHCA paper:
%   % ---------------------------------------------------------------------------
% Reference-Anchored Turn Taking (RATT)
%
% Standalone section intended to be merged into the main AHCA paper:
%   % ---------------------------------------------------------------------------
% Reference-Anchored Turn Taking (RATT)
%
% Standalone section intended to be merged into the main AHCA paper:
%   % ---------------------------------------------------------------------------
% Reference-Anchored Turn Taking (RATT)
%
% Standalone section intended to be merged into the main AHCA paper:
%   \input{reference_anchored_turn_taking}
% Requires only amsmath / amssymb, which the main document already loads.
% Every claim here maps to code in this repository:
%   lib/src/services/turn_taking.dart      (gate + lexical self-echo filter)
%   lib/src/services/call_session.dart     (turn-taking state machine)
%   android/.../FullDuplexAudioCapture.kt  (VOICE_COMMUNICATION + AEC attach)
% ---------------------------------------------------------------------------

\section{Reference-Anchored Turn Taking for Spoken Companion Agents}
\label{sec:ratt}

\subsection{Motivation}

A spoken companion agent differs from a conventional voice assistant in one
operational respect that turns out to be architecturally useful: it produces
its own audio. Every synthesised segment is generated locally, so the agent
holds the exact reference signal that a blind acoustic echo canceller has to
estimate. Existing pipelines either (i) mute the microphone during playback,
which forbids interruption and makes the agent feel like a broadcast, or
(ii) rely on the platform echo canceller alone, which is present on most
Android devices but whose residual return loss varies by more than 20\,dB
across hardware.

We exploit (i) being unnecessary: the known reference converts echo
suppression from \emph{blind system identification} into a
\emph{known-reference detection} problem, and (ii) being insufficient: the
hardware canceller removes the linear part of the echo, while a lexical filter
removes what the recogniser still hallucinates from the residual.

\subsection{Signal Model}

Let $f_s = 16$\,kHz and a frame of $T_f = 20$\,ms, giving $N = 320$ samples per
frame. At frame index $k$ the synthesiser reports the amplitude envelope
$a_k \in [0,1]$ of the segment currently being played (the per-frame RMS of the
generated WAV, normalised). The microphone observes

\begin{equation}
  x_k \;=\; \underbrace{\lambda\, a_{k-\delta}}_{\text{residual echo}}
        \;+\; \underbrace{s_k}_{\text{near-end speech}}
        \;+\; \underbrace{n_k}_{\text{noise}},
  \label{eq:observation}
\end{equation}

where $\lambda$ is the residual echo gain left by the echo canceller and
$\delta$ the bulk delay between playback and capture.

Two properties of $\delta$ matter. It is device dependent
(typically $20$--$120$\,ms) and it drifts with audio routing changes
(speaker, headset, Bluetooth). Estimating it online would require
cross-correlation over a window comparable to the reverberation time, which is
both expensive and fragile during double talk. Instead we discard the delay
entirely by replacing $a_{k-\delta}$ with a causal peak-hold of the reference:

\begin{equation}
  R_k \;=\; \max_{\,j \in [k-W,\,k]} a_j ,
  \qquad W = \left\lceil \frac{180\,\text{ms}}{T_f} \right\rceil = 9 .
  \label{eq:peakhold}
\end{equation}

Because $|a_j| \le 1$ and the envelope is piecewise smooth on the scale of a
syllable, $R_k$ is a tight upper bound on $a_{k-\delta}$ for every
$\delta \le W T_f$. The bound costs at most one frame of sensitivity, which is
irrelevant against the $H=140$\,ms evidence window used below.

\subsection{Decision Rule}

The predicted residual echo and the interruption threshold are

\begin{align}
  E_k &= \lambda R_k, \label{eq:echo}\\
  \tau_k &= \max\Big( \max\big(\tau_{\min},\; g\, \hat n_k\big),\;
             \gamma E_k + m \Big), \label{eq:threshold}
\end{align}

where $\hat n_k$ is a slow noise-floor estimate updated only when the agent is
silent, $g = 3.2$, $\tau_{\min} = 0.014$, $\gamma = 1.6$, and
$m = 0.012$.

Frame $k$ is \emph{voiced-evidence} when $x_k > \tau_k$. A barge-in is declared
only after voiced evidence has persisted for

\begin{equation}
  H = 140\,\text{ms} \quad (7\ \text{frames}),
\end{equation}

which rejects key clicks, door slams and single-syllable noise bursts while
keeping interruption latency under the ~200\,ms that users perceive as
simultaneous.

The leak gain is set from the measured capability of the platform, not from a
constant: $\lambda = 0.18$ when the echo canceller reported
\texttt{enabled}, and $\lambda = 0.55$ otherwise. The two regimes are not
interchangeable, and the agent never pretends they are (Section~\ref{sec:ratt-degrade}).

\subsection{Lexical Self-Echo Rejection}

Energy evidence answers \emph{when} to yield, not \emph{whether} the transcript
that follows belongs to the user. A recogniser fed a strong residual can emit
the agent's own sentence, which the naive pipeline then sends back to the
language model as a new user turn: a self-sustaining loop.

The agent again has an unfair advantage: it knows the exact text it read out.
Let $\sigma(\cdot)$ strip whitespace, punctuation and symbols, and let
$\mathcal{B}(u)$ be the multiset of character bigrams of $u$. For transcript
$u$ and a segment $v$ spoken within the last $8$\,s,

\begin{equation}
  D(u,v) \;=\; \frac{2 \sum_{b} \min\!\big(c_u(b), c_v(b)\big)}
                     {|\mathcal{B}(u)| + |\mathcal{B}(v)|},
  \label{eq:dice}
\end{equation}

is the bigram Dice coefficient. The transcript is rejected when
$D(u,v) \ge 0.6$, or when it is a substring of $v$ and
$|\sigma(u)| \ge 8$. Character bigrams are used rather than word tokens because
Chinese has no whitespace boundaries and ASR substitutions are typically
single-character. Transcripts shorter than six normalised characters are never
rejected, so a genuine short answer such as ``好的'' survives even if the agent
just said the same thing.

\subsection{Capability-Gated Degradation}
\label{sec:ratt-degrade}

The session queries the platform at capture start and records whether
\texttt{AcousticEchoCanceler} existed \emph{and} reported \texttt{enabled}
after attachment to the real \texttt{AudioRecord} session id. This is a
measurement, not a feature flag.

If the measurement fails, the session runs half duplex: the microphone is
closed before synthesis starts and reopened after playback plus a settle delay.
The UI states which mode is active, so an operator can tell a barge-in failure
from a device that never supported it. We consider this distinction a
correctness requirement rather than a product detail: silently degrading to
half duplex while claiming full duplex is the failure mode most voice-agent
papers do not report.

\subsection{Turn-Taking State Machine}

The session is a five-state machine
$\{\text{idle}, \text{listening}, \text{thinking}, \text{speaking},
\text{halted}\}$ in which every async continuation is guarded by a monotonically
increasing generation counter. This removes the class of bugs in which a late
completion from turn $t-1$ starts playback during turn $t$.

Two guards are specific to full duplex:

\begin{enumerate}
  \item \textbf{Post-dispatch grace.} For $900$\,ms after a turn is dispatched,
        barge-in from energy evidence is ignored. The tail of the utterance
        that opened the turn is still in the capture buffer; without the grace
        period it cancels the reply it just triggered.
  \item \textbf{Interrupt during generation.} In the \text{thinking} state the
        reference is zero, so \eqref{eq:threshold} collapses to the noise-floor
        term and any sustained speech is genuine. A user can therefore cut off
        the model mid-generation, not only mid-playback.
\end{enumerate}

\subsection{Complexity}

Per 20\,ms frame the gate performs $O(W)$ work for the peak-hold and $O(1)$
otherwise, with $W=9$: a sliding-window maximum maintained by a monotonic deque
gives amortised $O(1)$. The lexical filter runs once per final transcript over
at most four cached segments of $O(|v|)$ bigrams. Neither component allocates
inside the audio callback, and neither runs a timer: the gate is purely
event-driven by capture frames, so an idle session consumes no CPU.

\subsection{Evaluation Protocol}

We deliberately separate what is implemented from what remains to be measured.
The protocol below is what a reviewer should ask for; no numbers are claimed in
this section until they are produced by the harness.

\paragraph{Metrics}
\begin{itemize}
  \item \textbf{Barge-in latency}: time from the first user frame above
        $\tau_k$ to the last emitted playback sample. Target $<250$\,ms.
  \item \textbf{False barge-in rate}: interruptions with no true near-end
        speech, per hour of playback.
  \item \textbf{Self-echo turn rate}: assistant-voiced transcripts that reach
        the language model, with and without the lexical filter.
  \item \textbf{Over-suppression rate}: user utterances dropped by the
        filter, stratified by length.
  \item \textbf{Turn completion rate}: calls in which the user's utterance is
        followed by a spoken reply and a re-opened microphone.
\end{itemize}

\paragraph{Protocol}
Record the capture stream and the reference envelope simultaneously on the
target device. For each session, replay in four configurations:
(a) agent only, (b) user only, (c) agent and user overlapping (scripted
double talk at controlled signal-to-echo ratios), and (d) silence with room
noise. Configuration (a) bounds the false barge-in rate, (c) bounds latency and
missed interruptions, (d) bounds noise-triggered turns.

\paragraph{Ablations}
Peak-hold window $W \in \{5, 9, 15\}$; evidence window
$H \in \{80, 140, 220\}$\,ms; leak $\lambda \in \{0.12, 0.18, 0.35, 0.55\}$;
lexical filter on/off; hardware canceller on/off. Report $\lambda$ and $H$
jointly, since they trade latency against false interruption.

\paragraph{Device matrix}
At minimum one device per SoC family with a verified canceller and one without,
because the capability gate is the variable that changes the architecture
rather than a tuning parameter.

\subsection{Threats to Validity}

The peak-hold bound assumes the residual echo is a monotonically scaled copy of
the reference envelope; non-linear loudspeaker distortion at high volume
violates it, producing a higher effective $\lambda$. This is why the gate is
combined with the lexical filter rather than trusted alone. The lexical filter
cannot distinguish a user quoting the agent from the recogniser echoing it;
the $8$-character substring rule biases toward keeping the user's turn, which is
the correct error direction for an assistant. Finally, all thresholds were
chosen from the signal model and the frame budget, not fitted to a corpus; the
protocol above exists to test whether that transfer holds.

% Requires only amsmath / amssymb, which the main document already loads.
% Every claim here maps to code in this repository:
%   lib/src/services/turn_taking.dart      (gate + lexical self-echo filter)
%   lib/src/services/call_session.dart     (turn-taking state machine)
%   android/.../FullDuplexAudioCapture.kt  (VOICE_COMMUNICATION + AEC attach)
% ---------------------------------------------------------------------------

\section{Reference-Anchored Turn Taking for Spoken Companion Agents}
\label{sec:ratt}

\subsection{Motivation}

A spoken companion agent differs from a conventional voice assistant in one
operational respect that turns out to be architecturally useful: it produces
its own audio. Every synthesised segment is generated locally, so the agent
holds the exact reference signal that a blind acoustic echo canceller has to
estimate. Existing pipelines either (i) mute the microphone during playback,
which forbids interruption and makes the agent feel like a broadcast, or
(ii) rely on the platform echo canceller alone, which is present on most
Android devices but whose residual return loss varies by more than 20\,dB
across hardware.

We exploit (i) being unnecessary: the known reference converts echo
suppression from \emph{blind system identification} into a
\emph{known-reference detection} problem, and (ii) being insufficient: the
hardware canceller removes the linear part of the echo, while a lexical filter
removes what the recogniser still hallucinates from the residual.

\subsection{Signal Model}

Let $f_s = 16$\,kHz and a frame of $T_f = 20$\,ms, giving $N = 320$ samples per
frame. At frame index $k$ the synthesiser reports the amplitude envelope
$a_k \in [0,1]$ of the segment currently being played (the per-frame RMS of the
generated WAV, normalised). The microphone observes

\begin{equation}
  x_k \;=\; \underbrace{\lambda\, a_{k-\delta}}_{\text{residual echo}}
        \;+\; \underbrace{s_k}_{\text{near-end speech}}
        \;+\; \underbrace{n_k}_{\text{noise}},
  \label{eq:observation}
\end{equation}

where $\lambda$ is the residual echo gain left by the echo canceller and
$\delta$ the bulk delay between playback and capture.

Two properties of $\delta$ matter. It is device dependent
(typically $20$--$120$\,ms) and it drifts with audio routing changes
(speaker, headset, Bluetooth). Estimating it online would require
cross-correlation over a window comparable to the reverberation time, which is
both expensive and fragile during double talk. Instead we discard the delay
entirely by replacing $a_{k-\delta}$ with a causal peak-hold of the reference:

\begin{equation}
  R_k \;=\; \max_{\,j \in [k-W,\,k]} a_j ,
  \qquad W = \left\lceil \frac{180\,\text{ms}}{T_f} \right\rceil = 9 .
  \label{eq:peakhold}
\end{equation}

Because $|a_j| \le 1$ and the envelope is piecewise smooth on the scale of a
syllable, $R_k$ is a tight upper bound on $a_{k-\delta}$ for every
$\delta \le W T_f$. The bound costs at most one frame of sensitivity, which is
irrelevant against the $H=140$\,ms evidence window used below.

\subsection{Decision Rule}

The predicted residual echo and the interruption threshold are

\begin{align}
  E_k &= \lambda R_k, \label{eq:echo}\\
  \tau_k &= \max\Big( \max\big(\tau_{\min},\; g\, \hat n_k\big),\;
             \gamma E_k + m \Big), \label{eq:threshold}
\end{align}

where $\hat n_k$ is a slow noise-floor estimate updated only when the agent is
silent, $g = 3.2$, $\tau_{\min} = 0.014$, $\gamma = 1.6$, and
$m = 0.012$.

Frame $k$ is \emph{voiced-evidence} when $x_k > \tau_k$. A barge-in is declared
only after voiced evidence has persisted for

\begin{equation}
  H = 140\,\text{ms} \quad (7\ \text{frames}),
\end{equation}

which rejects key clicks, door slams and single-syllable noise bursts while
keeping interruption latency under the ~200\,ms that users perceive as
simultaneous.

The leak gain is set from the measured capability of the platform, not from a
constant: $\lambda = 0.18$ when the echo canceller reported
\texttt{enabled}, and $\lambda = 0.55$ otherwise. The two regimes are not
interchangeable, and the agent never pretends they are (Section~\ref{sec:ratt-degrade}).

\subsection{Lexical Self-Echo Rejection}

Energy evidence answers \emph{when} to yield, not \emph{whether} the transcript
that follows belongs to the user. A recogniser fed a strong residual can emit
the agent's own sentence, which the naive pipeline then sends back to the
language model as a new user turn: a self-sustaining loop.

The agent again has an unfair advantage: it knows the exact text it read out.
Let $\sigma(\cdot)$ strip whitespace, punctuation and symbols, and let
$\mathcal{B}(u)$ be the multiset of character bigrams of $u$. For transcript
$u$ and a segment $v$ spoken within the last $8$\,s,

\begin{equation}
  D(u,v) \;=\; \frac{2 \sum_{b} \min\!\big(c_u(b), c_v(b)\big)}
                     {|\mathcal{B}(u)| + |\mathcal{B}(v)|},
  \label{eq:dice}
\end{equation}

is the bigram Dice coefficient. The transcript is rejected when
$D(u,v) \ge 0.6$, or when it is a substring of $v$ and
$|\sigma(u)| \ge 8$. Character bigrams are used rather than word tokens because
Chinese has no whitespace boundaries and ASR substitutions are typically
single-character. Transcripts shorter than six normalised characters are never
rejected, so a genuine short answer such as ``好的'' survives even if the agent
just said the same thing.

\subsection{Capability-Gated Degradation}
\label{sec:ratt-degrade}

The session queries the platform at capture start and records whether
\texttt{AcousticEchoCanceler} existed \emph{and} reported \texttt{enabled}
after attachment to the real \texttt{AudioRecord} session id. This is a
measurement, not a feature flag.

If the measurement fails, the session runs half duplex: the microphone is
closed before synthesis starts and reopened after playback plus a settle delay.
The UI states which mode is active, so an operator can tell a barge-in failure
from a device that never supported it. We consider this distinction a
correctness requirement rather than a product detail: silently degrading to
half duplex while claiming full duplex is the failure mode most voice-agent
papers do not report.

\subsection{Turn-Taking State Machine}

The session is a five-state machine
$\{\text{idle}, \text{listening}, \text{thinking}, \text{speaking},
\text{halted}\}$ in which every async continuation is guarded by a monotonically
increasing generation counter. This removes the class of bugs in which a late
completion from turn $t-1$ starts playback during turn $t$.

Two guards are specific to full duplex:

\begin{enumerate}
  \item \textbf{Post-dispatch grace.} For $900$\,ms after a turn is dispatched,
        barge-in from energy evidence is ignored. The tail of the utterance
        that opened the turn is still in the capture buffer; without the grace
        period it cancels the reply it just triggered.
  \item \textbf{Interrupt during generation.} In the \text{thinking} state the
        reference is zero, so \eqref{eq:threshold} collapses to the noise-floor
        term and any sustained speech is genuine. A user can therefore cut off
        the model mid-generation, not only mid-playback.
\end{enumerate}

\subsection{Complexity}

Per 20\,ms frame the gate performs $O(W)$ work for the peak-hold and $O(1)$
otherwise, with $W=9$: a sliding-window maximum maintained by a monotonic deque
gives amortised $O(1)$. The lexical filter runs once per final transcript over
at most four cached segments of $O(|v|)$ bigrams. Neither component allocates
inside the audio callback, and neither runs a timer: the gate is purely
event-driven by capture frames, so an idle session consumes no CPU.

\subsection{Evaluation Protocol}

We deliberately separate what is implemented from what remains to be measured.
The protocol below is what a reviewer should ask for; no numbers are claimed in
this section until they are produced by the harness.

\paragraph{Metrics}
\begin{itemize}
  \item \textbf{Barge-in latency}: time from the first user frame above
        $\tau_k$ to the last emitted playback sample. Target $<250$\,ms.
  \item \textbf{False barge-in rate}: interruptions with no true near-end
        speech, per hour of playback.
  \item \textbf{Self-echo turn rate}: assistant-voiced transcripts that reach
        the language model, with and without the lexical filter.
  \item \textbf{Over-suppression rate}: user utterances dropped by the
        filter, stratified by length.
  \item \textbf{Turn completion rate}: calls in which the user's utterance is
        followed by a spoken reply and a re-opened microphone.
\end{itemize}

\paragraph{Protocol}
Record the capture stream and the reference envelope simultaneously on the
target device. For each session, replay in four configurations:
(a) agent only, (b) user only, (c) agent and user overlapping (scripted
double talk at controlled signal-to-echo ratios), and (d) silence with room
noise. Configuration (a) bounds the false barge-in rate, (c) bounds latency and
missed interruptions, (d) bounds noise-triggered turns.

\paragraph{Ablations}
Peak-hold window $W \in \{5, 9, 15\}$; evidence window
$H \in \{80, 140, 220\}$\,ms; leak $\lambda \in \{0.12, 0.18, 0.35, 0.55\}$;
lexical filter on/off; hardware canceller on/off. Report $\lambda$ and $H$
jointly, since they trade latency against false interruption.

\paragraph{Device matrix}
At minimum one device per SoC family with a verified canceller and one without,
because the capability gate is the variable that changes the architecture
rather than a tuning parameter.

\subsection{Threats to Validity}

The peak-hold bound assumes the residual echo is a monotonically scaled copy of
the reference envelope; non-linear loudspeaker distortion at high volume
violates it, producing a higher effective $\lambda$. This is why the gate is
combined with the lexical filter rather than trusted alone. The lexical filter
cannot distinguish a user quoting the agent from the recogniser echoing it;
the $8$-character substring rule biases toward keeping the user's turn, which is
the correct error direction for an assistant. Finally, all thresholds were
chosen from the signal model and the frame budget, not fitted to a corpus; the
protocol above exists to test whether that transfer holds.

% Requires only amsmath / amssymb, which the main document already loads.
% Every claim here maps to code in this repository:
%   lib/src/services/turn_taking.dart      (gate + lexical self-echo filter)
%   lib/src/services/call_session.dart     (turn-taking state machine)
%   android/.../FullDuplexAudioCapture.kt  (VOICE_COMMUNICATION + AEC attach)
% ---------------------------------------------------------------------------

\section{Reference-Anchored Turn Taking for Spoken Companion Agents}
\label{sec:ratt}

\subsection{Motivation}

A spoken companion agent differs from a conventional voice assistant in one
operational respect that turns out to be architecturally useful: it produces
its own audio. Every synthesised segment is generated locally, so the agent
holds the exact reference signal that a blind acoustic echo canceller has to
estimate. Existing pipelines either (i) mute the microphone during playback,
which forbids interruption and makes the agent feel like a broadcast, or
(ii) rely on the platform echo canceller alone, which is present on most
Android devices but whose residual return loss varies by more than 20\,dB
across hardware.

We exploit (i) being unnecessary: the known reference converts echo
suppression from \emph{blind system identification} into a
\emph{known-reference detection} problem, and (ii) being insufficient: the
hardware canceller removes the linear part of the echo, while a lexical filter
removes what the recogniser still hallucinates from the residual.

\subsection{Signal Model}

Let $f_s = 16$\,kHz and a frame of $T_f = 20$\,ms, giving $N = 320$ samples per
frame. At frame index $k$ the synthesiser reports the amplitude envelope
$a_k \in [0,1]$ of the segment currently being played (the per-frame RMS of the
generated WAV, normalised). The microphone observes

\begin{equation}
  x_k \;=\; \underbrace{\lambda\, a_{k-\delta}}_{\text{residual echo}}
        \;+\; \underbrace{s_k}_{\text{near-end speech}}
        \;+\; \underbrace{n_k}_{\text{noise}},
  \label{eq:observation}
\end{equation}

where $\lambda$ is the residual echo gain left by the echo canceller and
$\delta$ the bulk delay between playback and capture.

Two properties of $\delta$ matter. It is device dependent
(typically $20$--$120$\,ms) and it drifts with audio routing changes
(speaker, headset, Bluetooth). Estimating it online would require
cross-correlation over a window comparable to the reverberation time, which is
both expensive and fragile during double talk. Instead we discard the delay
entirely by replacing $a_{k-\delta}$ with a causal peak-hold of the reference:

\begin{equation}
  R_k \;=\; \max_{\,j \in [k-W,\,k]} a_j ,
  \qquad W = \left\lceil \frac{180\,\text{ms}}{T_f} \right\rceil = 9 .
  \label{eq:peakhold}
\end{equation}

Because $|a_j| \le 1$ and the envelope is piecewise smooth on the scale of a
syllable, $R_k$ is a tight upper bound on $a_{k-\delta}$ for every
$\delta \le W T_f$. The bound costs at most one frame of sensitivity, which is
irrelevant against the $H=140$\,ms evidence window used below.

\subsection{Decision Rule}

The predicted residual echo and the interruption threshold are

\begin{align}
  E_k &= \lambda R_k, \label{eq:echo}\\
  \tau_k &= \max\Big( \max\big(\tau_{\min},\; g\, \hat n_k\big),\;
             \gamma E_k + m \Big), \label{eq:threshold}
\end{align}

where $\hat n_k$ is a slow noise-floor estimate updated only when the agent is
silent, $g = 3.2$, $\tau_{\min} = 0.014$, $\gamma = 1.6$, and
$m = 0.012$.

Frame $k$ is \emph{voiced-evidence} when $x_k > \tau_k$. A barge-in is declared
only after voiced evidence has persisted for

\begin{equation}
  H = 140\,\text{ms} \quad (7\ \text{frames}),
\end{equation}

which rejects key clicks, door slams and single-syllable noise bursts while
keeping interruption latency under the ~200\,ms that users perceive as
simultaneous.

The leak gain is set from the measured capability of the platform, not from a
constant: $\lambda = 0.18$ when the echo canceller reported
\texttt{enabled}, and $\lambda = 0.55$ otherwise. The two regimes are not
interchangeable, and the agent never pretends they are (Section~\ref{sec:ratt-degrade}).

\subsection{Lexical Self-Echo Rejection}

Energy evidence answers \emph{when} to yield, not \emph{whether} the transcript
that follows belongs to the user. A recogniser fed a strong residual can emit
the agent's own sentence, which the naive pipeline then sends back to the
language model as a new user turn: a self-sustaining loop.

The agent again has an unfair advantage: it knows the exact text it read out.
Let $\sigma(\cdot)$ strip whitespace, punctuation and symbols, and let
$\mathcal{B}(u)$ be the multiset of character bigrams of $u$. For transcript
$u$ and a segment $v$ spoken within the last $8$\,s,

\begin{equation}
  D(u,v) \;=\; \frac{2 \sum_{b} \min\!\big(c_u(b), c_v(b)\big)}
                     {|\mathcal{B}(u)| + |\mathcal{B}(v)|},
  \label{eq:dice}
\end{equation}

is the bigram Dice coefficient. The transcript is rejected when
$D(u,v) \ge 0.6$, or when it is a substring of $v$ and
$|\sigma(u)| \ge 8$. Character bigrams are used rather than word tokens because
Chinese has no whitespace boundaries and ASR substitutions are typically
single-character. Transcripts shorter than six normalised characters are never
rejected, so a genuine short answer such as ``好的'' survives even if the agent
just said the same thing.

\subsection{Capability-Gated Degradation}
\label{sec:ratt-degrade}

The session queries the platform at capture start and records whether
\texttt{AcousticEchoCanceler} existed \emph{and} reported \texttt{enabled}
after attachment to the real \texttt{AudioRecord} session id. This is a
measurement, not a feature flag.

If the measurement fails, the session runs half duplex: the microphone is
closed before synthesis starts and reopened after playback plus a settle delay.
The UI states which mode is active, so an operator can tell a barge-in failure
from a device that never supported it. We consider this distinction a
correctness requirement rather than a product detail: silently degrading to
half duplex while claiming full duplex is the failure mode most voice-agent
papers do not report.

\subsection{Turn-Taking State Machine}

The session is a five-state machine
$\{\text{idle}, \text{listening}, \text{thinking}, \text{speaking},
\text{halted}\}$ in which every async continuation is guarded by a monotonically
increasing generation counter. This removes the class of bugs in which a late
completion from turn $t-1$ starts playback during turn $t$.

Two guards are specific to full duplex:

\begin{enumerate}
  \item \textbf{Post-dispatch grace.} For $900$\,ms after a turn is dispatched,
        barge-in from energy evidence is ignored. The tail of the utterance
        that opened the turn is still in the capture buffer; without the grace
        period it cancels the reply it just triggered.
  \item \textbf{Interrupt during generation.} In the \text{thinking} state the
        reference is zero, so \eqref{eq:threshold} collapses to the noise-floor
        term and any sustained speech is genuine. A user can therefore cut off
        the model mid-generation, not only mid-playback.
\end{enumerate}

\subsection{Complexity}

Per 20\,ms frame the gate performs $O(W)$ work for the peak-hold and $O(1)$
otherwise, with $W=9$: a sliding-window maximum maintained by a monotonic deque
gives amortised $O(1)$. The lexical filter runs once per final transcript over
at most four cached segments of $O(|v|)$ bigrams. Neither component allocates
inside the audio callback, and neither runs a timer: the gate is purely
event-driven by capture frames, so an idle session consumes no CPU.

\subsection{Evaluation Protocol}

We deliberately separate what is implemented from what remains to be measured.
The protocol below is what a reviewer should ask for; no numbers are claimed in
this section until they are produced by the harness.

\paragraph{Metrics}
\begin{itemize}
  \item \textbf{Barge-in latency}: time from the first user frame above
        $\tau_k$ to the last emitted playback sample. Target $<250$\,ms.
  \item \textbf{False barge-in rate}: interruptions with no true near-end
        speech, per hour of playback.
  \item \textbf{Self-echo turn rate}: assistant-voiced transcripts that reach
        the language model, with and without the lexical filter.
  \item \textbf{Over-suppression rate}: user utterances dropped by the
        filter, stratified by length.
  \item \textbf{Turn completion rate}: calls in which the user's utterance is
        followed by a spoken reply and a re-opened microphone.
\end{itemize}

\paragraph{Protocol}
Record the capture stream and the reference envelope simultaneously on the
target device. For each session, replay in four configurations:
(a) agent only, (b) user only, (c) agent and user overlapping (scripted
double talk at controlled signal-to-echo ratios), and (d) silence with room
noise. Configuration (a) bounds the false barge-in rate, (c) bounds latency and
missed interruptions, (d) bounds noise-triggered turns.

\paragraph{Ablations}
Peak-hold window $W \in \{5, 9, 15\}$; evidence window
$H \in \{80, 140, 220\}$\,ms; leak $\lambda \in \{0.12, 0.18, 0.35, 0.55\}$;
lexical filter on/off; hardware canceller on/off. Report $\lambda$ and $H$
jointly, since they trade latency against false interruption.

\paragraph{Device matrix}
At minimum one device per SoC family with a verified canceller and one without,
because the capability gate is the variable that changes the architecture
rather than a tuning parameter.

\subsection{Threats to Validity}

The peak-hold bound assumes the residual echo is a monotonically scaled copy of
the reference envelope; non-linear loudspeaker distortion at high volume
violates it, producing a higher effective $\lambda$. This is why the gate is
combined with the lexical filter rather than trusted alone. The lexical filter
cannot distinguish a user quoting the agent from the recogniser echoing it;
the $8$-character substring rule biases toward keeping the user's turn, which is
the correct error direction for an assistant. Finally, all thresholds were
chosen from the signal model and the frame budget, not fitted to a corpus; the
protocol above exists to test whether that transfer holds.

% Requires only amsmath / amssymb, which the main document already loads.
% Every claim here maps to code in this repository:
%   lib/src/services/turn_taking.dart      (gate + lexical self-echo filter)
%   lib/src/services/call_session.dart     (turn-taking state machine)
%   android/.../FullDuplexAudioCapture.kt  (VOICE_COMMUNICATION + AEC attach)
% ---------------------------------------------------------------------------

\section{Reference-Anchored Turn Taking for Spoken Companion Agents}
\label{sec:ratt}

\subsection{Motivation}

A spoken companion agent differs from a conventional voice assistant in one
operational respect that turns out to be architecturally useful: it produces
its own audio. Every synthesised segment is generated locally, so the agent
holds the exact reference signal that a blind acoustic echo canceller has to
estimate. Existing pipelines either (i) mute the microphone during playback,
which forbids interruption and makes the agent feel like a broadcast, or
(ii) rely on the platform echo canceller alone, which is present on most
Android devices but whose residual return loss varies by more than 20\,dB
across hardware.

We exploit (i) being unnecessary: the known reference converts echo
suppression from \emph{blind system identification} into a
\emph{known-reference detection} problem, and (ii) being insufficient: the
hardware canceller removes the linear part of the echo, while a lexical filter
removes what the recogniser still hallucinates from the residual.

\subsection{Signal Model}

Let $f_s = 16$\,kHz and a frame of $T_f = 20$\,ms, giving $N = 320$ samples per
frame. At frame index $k$ the synthesiser reports the amplitude envelope
$a_k \in [0,1]$ of the segment currently being played (the per-frame RMS of the
generated WAV, normalised). The microphone observes

\begin{equation}
  x_k \;=\; \underbrace{\lambda\, a_{k-\delta}}_{\text{residual echo}}
        \;+\; \underbrace{s_k}_{\text{near-end speech}}
        \;+\; \underbrace{n_k}_{\text{noise}},
  \label{eq:observation}
\end{equation}

where $\lambda$ is the residual echo gain left by the echo canceller and
$\delta$ the bulk delay between playback and capture.

Two properties of $\delta$ matter. It is device dependent
(typically $20$--$120$\,ms) and it drifts with audio routing changes
(speaker, headset, Bluetooth). Estimating it online would require
cross-correlation over a window comparable to the reverberation time, which is
both expensive and fragile during double talk. Instead we discard the delay
entirely by replacing $a_{k-\delta}$ with a causal peak-hold of the reference:

\begin{equation}
  R_k \;=\; \max_{\,j \in [k-W,\,k]} a_j ,
  \qquad W = \left\lceil \frac{180\,\text{ms}}{T_f} \right\rceil = 9 .
  \label{eq:peakhold}
\end{equation}

Because $|a_j| \le 1$ and the envelope is piecewise smooth on the scale of a
syllable, $R_k$ is a tight upper bound on $a_{k-\delta}$ for every
$\delta \le W T_f$. The bound costs at most one frame of sensitivity, which is
irrelevant against the $H=140$\,ms evidence window used below.

\subsection{Decision Rule}

The predicted residual echo and the interruption threshold are

\begin{align}
  E_k &= \lambda R_k, \label{eq:echo}\\
  \tau_k &= \max\Big( \max\big(\tau_{\min},\; g\, \hat n_k\big),\;
             \gamma E_k + m \Big), \label{eq:threshold}
\end{align}

where $\hat n_k$ is a slow noise-floor estimate updated only when the agent is
silent, $g = 3.2$, $\tau_{\min} = 0.014$, $\gamma = 1.6$, and
$m = 0.012$.

Frame $k$ is \emph{voiced-evidence} when $x_k > \tau_k$. A barge-in is declared
only after voiced evidence has persisted for

\begin{equation}
  H = 140\,\text{ms} \quad (7\ \text{frames}),
\end{equation}

which rejects key clicks, door slams and single-syllable noise bursts while
keeping interruption latency under the ~200\,ms that users perceive as
simultaneous.

The leak gain is set from the measured capability of the platform, not from a
constant: $\lambda = 0.18$ when the echo canceller reported
\texttt{enabled}, and $\lambda = 0.55$ otherwise. The two regimes are not
interchangeable, and the agent never pretends they are (Section~\ref{sec:ratt-degrade}).

\subsection{Lexical Self-Echo Rejection}

Energy evidence answers \emph{when} to yield, not \emph{whether} the transcript
that follows belongs to the user. A recogniser fed a strong residual can emit
the agent's own sentence, which the naive pipeline then sends back to the
language model as a new user turn: a self-sustaining loop.

The agent again has an unfair advantage: it knows the exact text it read out.
Let $\sigma(\cdot)$ strip whitespace, punctuation and symbols, and let
$\mathcal{B}(u)$ be the multiset of character bigrams of $u$. For transcript
$u$ and a segment $v$ spoken within the last $8$\,s,

\begin{equation}
  D(u,v) \;=\; \frac{2 \sum_{b} \min\!\big(c_u(b), c_v(b)\big)}
                     {|\mathcal{B}(u)| + |\mathcal{B}(v)|},
  \label{eq:dice}
\end{equation}

is the bigram Dice coefficient. The transcript is rejected when
$D(u,v) \ge 0.6$, or when it is a substring of $v$ and
$|\sigma(u)| \ge 8$. Character bigrams are used rather than word tokens because
Chinese has no whitespace boundaries and ASR substitutions are typically
single-character. Transcripts shorter than six normalised characters are never
rejected, so a genuine short answer such as ``好的'' survives even if the agent
just said the same thing.

\subsection{Capability-Gated Degradation}
\label{sec:ratt-degrade}

The session queries the platform at capture start and records whether
\texttt{AcousticEchoCanceler} existed \emph{and} reported \texttt{enabled}
after attachment to the real \texttt{AudioRecord} session id. This is a
measurement, not a feature flag.

If the measurement fails, the session runs half duplex: the microphone is
closed before synthesis starts and reopened after playback plus a settle delay.
The UI states which mode is active, so an operator can tell a barge-in failure
from a device that never supported it. We consider this distinction a
correctness requirement rather than a product detail: silently degrading to
half duplex while claiming full duplex is the failure mode most voice-agent
papers do not report.

\subsection{Turn-Taking State Machine}

The session is a five-state machine
$\{\text{idle}, \text{listening}, \text{thinking}, \text{speaking},
\text{halted}\}$ in which every async continuation is guarded by a monotonically
increasing generation counter. This removes the class of bugs in which a late
completion from turn $t-1$ starts playback during turn $t$.

Two guards are specific to full duplex:

\begin{enumerate}
  \item \textbf{Post-dispatch grace.} For $900$\,ms after a turn is dispatched,
        barge-in from energy evidence is ignored. The tail of the utterance
        that opened the turn is still in the capture buffer; without the grace
        period it cancels the reply it just triggered.
  \item \textbf{Interrupt during generation.} In the \text{thinking} state the
        reference is zero, so \eqref{eq:threshold} collapses to the noise-floor
        term and any sustained speech is genuine. A user can therefore cut off
        the model mid-generation, not only mid-playback.
\end{enumerate}

\subsection{Complexity}

Per 20\,ms frame the gate performs $O(W)$ work for the peak-hold and $O(1)$
otherwise, with $W=9$: a sliding-window maximum maintained by a monotonic deque
gives amortised $O(1)$. The lexical filter runs once per final transcript over
at most four cached segments of $O(|v|)$ bigrams. Neither component allocates
inside the audio callback, and neither runs a timer: the gate is purely
event-driven by capture frames, so an idle session consumes no CPU.

\subsection{Evaluation Protocol}

We deliberately separate what is implemented from what remains to be measured.
The protocol below is what a reviewer should ask for; no numbers are claimed in
this section until they are produced by the harness.

\paragraph{Metrics}
\begin{itemize}
  \item \textbf{Barge-in latency}: time from the first user frame above
        $\tau_k$ to the last emitted playback sample. Target $<250$\,ms.
  \item \textbf{False barge-in rate}: interruptions with no true near-end
        speech, per hour of playback.
  \item \textbf{Self-echo turn rate}: assistant-voiced transcripts that reach
        the language model, with and without the lexical filter.
  \item \textbf{Over-suppression rate}: user utterances dropped by the
        filter, stratified by length.
  \item \textbf{Turn completion rate}: calls in which the user's utterance is
        followed by a spoken reply and a re-opened microphone.
\end{itemize}

\paragraph{Protocol}
Record the capture stream and the reference envelope simultaneously on the
target device. For each session, replay in four configurations:
(a) agent only, (b) user only, (c) agent and user overlapping (scripted
double talk at controlled signal-to-echo ratios), and (d) silence with room
noise. Configuration (a) bounds the false barge-in rate, (c) bounds latency and
missed interruptions, (d) bounds noise-triggered turns.

\paragraph{Ablations}
Peak-hold window $W \in \{5, 9, 15\}$; evidence window
$H \in \{80, 140, 220\}$\,ms; leak $\lambda \in \{0.12, 0.18, 0.35, 0.55\}$;
lexical filter on/off; hardware canceller on/off. Report $\lambda$ and $H$
jointly, since they trade latency against false interruption.

\paragraph{Device matrix}
At minimum one device per SoC family with a verified canceller and one without,
because the capability gate is the variable that changes the architecture
rather than a tuning parameter.

\subsection{Threats to Validity}

The peak-hold bound assumes the residual echo is a monotonically scaled copy of
the reference envelope; non-linear loudspeaker distortion at high volume
violates it, producing a higher effective $\lambda$. This is why the gate is
combined with the lexical filter rather than trusted alone. The lexical filter
cannot distinguish a user quoting the agent from the recogniser echoing it;
the $8$-character substring rule biases toward keeping the user's turn, which is
the correct error direction for an assistant. Finally, all thresholds were
chosen from the signal model and the frame budget, not fitted to a corpus; the
protocol above exists to test whether that transfer holds.
